\chapter{Materia a partir de los alimentos}\label{ch:g6:matter-from-food}

Un ternero de \qty{40}{kg} se convierte en dos años en una vaca de
\qty{600}{kg}. Media tonelada de vaca nueva --- carne, \omterm{def:g3:bones-and-muscles:skeleton}{hueso}, cuero,
cuerno --- apareció de alguna parte. La vaca comió hierba; la materia
nueva \emph{no es hierba}; y sin embargo, sin la hierba no existiría
nada de ella. Lo que hacen los \omterm{def:g1:animals-around-us:animal}{animales} con su comida es el asunto de
este capítulo --- el lado de los \omterm{def:g5:food-webs:roles}{consumidores} en la contabilidad del
\cref{ch:g6:plants-are-producers}.

\section{Los animales construyen su materia con los alimentos}

\begin{proposition}[Los consumidores transforman materia orgánica]\label{prop:g6:matter-from-food:transform}
Un \omterm{def:g1:animals-around-us:animal}{animal} construye su propia \omterm{def:g6:plants-are-producers:matter}{materia orgánica} --- \omterm{def:g3:bones-and-muscles:muscle}{músculo}, \omterm{def:g3:bones-and-muscles:skeleton}{hueso},
grasa, \omterm{def:g1:animals-around-us:covering}{pluma}, \omterm{def:g1:animals-around-us:covering}{concha} --- a partir de la \emph{\omterm{def:g6:plants-are-producers:matter}{materia orgánica} que
come}. La \omterm{def:g4:journey-of-food:digestion}{digestión} (\cref{def:g4:journey-of-food:digestion}) desmonta
el alimento en \omterm{def:g4:journey-of-food:digestion}{nutrientes}; la \omterm{prop:g4:journey-of-food:absorb}{sangre} los reparte
(\cref{prop:g4:heart-and-blood:carries}); y las \omterm{def:g5:first-look-at-cells:cell}{células} del cuerpo
(\cref{prop:g5:first-look-at-cells:growth}) los vuelven a montar como
sustancia propia del \omterm{def:g1:animals-around-us:animal}{animal}. Los \omterm{def:g1:animals-around-us:animal}{animales} son \emph{transformadores} de
materia, nunca fabricantes a partir de lo mineral: ese oficio es solo de
las \omterm{def:g5:food-webs:roles}{productoras}.
\end{proposition}
\admitted

\begin{example}[De hierba a vaca]\label{ex:g6:matter-from-food:cow}
Sigue la media tonelada: la vaca pasta hora tras hora (la
\cref{prop:g2:what-animals-eat:daily} explicó esas horas tan largas);
su pesada \omterm{def:g4:journey-of-food:digestion}{digestión} (\cref{exo:g4:journey-of-food:11}) desmonta la
hierba en \omterm{def:g4:journey-of-food:digestion}{nutrientes}; su \omterm{prop:g4:journey-of-food:absorb}{sangre} los lleva a la ubre, al \omterm{def:g3:bones-and-muscles:muscle}{músculo} y al
ternero que crece dentro de ella; y sus \omterm{def:g5:first-look-at-cells:cell}{células} construyen materia de
vaca a partir de materia de hierba. El pasto está, literalmente,
paseándose por ahí con otra forma.
\end{example}

\begin{remark}[La materia cambia de manos, no de naturaleza]\label{rem:g6:matter-from-food:hands}
En cada flecha de una \omterm{def:g3:food-chains:chain}{cadena alimentaria}
(\cref{def:g3:food-chains:chain}), la \omterm{def:g6:plants-are-producers:matter}{materia orgánica} cambia de dueño y
se reconstruye --- materia de hierba en materia de vaca, materia de
vaca (la \omterm{def:g2:animals-grow-up:born}{leche}, la carne) en materia de niño. Lo \emph{orgánico} no sale
nunca de la nada: los \omterm{def:g5:food-webs:roles}{consumidores} se lo pasan y le dan otra forma. Solo
se fabricó de nuevo en el arranque verde de la cadena
(\cref{prop:g6:plants-are-producers:firstlink}); y solo en el final de
los \omterm{def:g5:food-webs:roles}{descomponedores}, en el capítulo siguiente, vuelve a ser mineral.
\end{remark}

\section{Crecer es producir materia}

\begin{definition}[La producción de materia de un animal]\label{def:g6:matter-from-food:production}
La \emph{producción de materia}\index{producción de materia} de un
\omterm{def:g1:animals-around-us:animal}{animal} es la \omterm{def:g6:plants-are-producers:matter}{materia orgánica} nueva que construye: el crecimiento de las
crías (los kilos que gana el ternero), las renovaciones anuales (el \omterm{def:g1:animals-around-us:covering}{pelo}
que se muda y vuelve a salir, las astas, las \omterm{def:g1:animals-around-us:covering}{plumas}), las reparaciones
(el \omterm{def:g3:bones-and-muscles:skeleton}{hueso} soldado de la \cref{rem:g3:bones-and-muscles:alive}) y los
productos (\omterm{def:g2:animals-grow-up:born}{leche}, \omterm{def:g2:animals-grow-up:born}{huevos}, lana). Todo ello es alimento reconstruido.
\end{definition}

\begin{example}[Pesar la producción]\label{ex:g6:matter-from-food:weighing}
Los ganaderos pesan la producción a diario: una vaca lechera da unos
\qty{25}{kg} de \omterm{def:g2:animals-grow-up:born}{leche} al día --- \omterm{def:g6:plants-are-producers:matter}{materia orgánica} que tiene que
reconstruir con el pienso, día tras día; el \omterm{def:g2:animals-grow-up:born}{huevo} de una gallina son
unos \qty{60}{g} de provisiones empaquetadas
(\cref{rem:g2:animals-grow-up:egg}), montadas de la noche a la mañana
con su grano. Sin pienso no hay producción: la contabilidad cuadra sin
piedad.
\end{example}

\begin{proposition}[No todo el alimento se convierte en cuerpo]\label{prop:g6:matter-from-food:losses}
Un \omterm{def:g5:food-webs:roles}{consumidor} no convierte nunca todo su alimento en materia nueva:
\begin{itemize}
  \item una parte no llega a absorberse y sale en forma de excrementos
        (\cref{prop:g4:journey-of-food:absorb});
  \item una parte grande se \emph{quema} --- combinada con el oxígeno
        de la \omterm{def:g4:breathing:movements}{respiración} (\cref{prop:g4:breathing:exchange}) --- para
        dar fuerza al movimiento y calor: materia gastada como
        combustible, que sale en forma de gas espirado;
  \item solo el resto se convierte en crecimiento y en productos.
\end{itemize}
De la hierba que se come una vaca, solo una fracción modesta acaba
siendo vaca nueva --- que es por lo que cada piso de la pirámide de la
\cref{prop:g5:food-webs:pyramid} es mucho más pequeño que el de debajo.
\end{proposition}
\admitted

\begin{example}[La pirámide, explicada por fin]\label{ex:g6:matter-from-food:pyramid}
Las toneladas de hierba de un prado alimentan a una manada de conejos;
los kilos de los conejos alimentan a un zorro o dos. En cada flecha,
casi toda la materia se quema para vivir o se pierde, y solo una
fracción se pasa hacia arriba como cuerpo. Los grandes cazadores son
escasos (\cref{prop:g5:food-webs:pyramid}) porque viven de los
\emph{restos de los restos} --- es aritmética, no mala suerte.
\end{example}

\section{Menús, dientes y un mismo oficio}

\begin{proposition}[Menús distintos, la misma transformación]\label{prop:g6:matter-from-food:menus}
\omterm{def:g2:what-animals-eat:kinds}{Herbívoro}, \omterm{def:g2:what-animals-eat:kinds}{carnívoro} y \omterm{def:g2:what-animals-eat:kinds}{omnívoro} (\cref{def:g2:what-animals-eat:kinds})
nombran \emph{fuentes} distintas de \omterm{def:g6:plants-are-producers:matter}{materia orgánica}, no oficios
distintos. La vaca que reconstruye hierba y el zorro que reconstruye
conejo hacen los dos la transformación de la
\cref{prop:g6:matter-from-food:transform} --- con el equipo ajustado al
menú (\cref{prop:g2:what-animals-eat:match},
\cref{exo:g4:journey-of-food:11}): la hierba exige muelas y una
\omterm{def:g4:journey-of-food:digestion}{digestión} enorme, la carne recompensa a los dientes de agarrar y a un
\omterm{def:g4:journey-of-food:tract}{tubo digestivo} corto.
\end{proposition}
\admitted

\begin{method}[Auditar la materia de un animal]\label{met:g6:matter-from-food:audit}
Para cualquier \omterm{def:g1:animals-around-us:animal}{animal}, monta su contabilidad de materia:
\begin{enumerate}
  \item \emph{entradas}: ¿qué \omterm{def:g6:plants-are-producers:matter}{materia orgánica} come y de quién? (sus
        flechas de entrada en la \omterm{def:g5:food-webs:web}{red alimentaria},
        \cref{def:g5:food-webs:web});
  \item \emph{construido}: ¿qué crecimiento, renovaciones y productos
        fabrica?
  \item \emph{salidas}: excrementos y materia quemada para el
        movimiento y el calor;
  \item comprueba el balance: la producción tiene que ser el resto
        modesto --- si una afirmación la hace mayor (``esta charca da
        más materia de \omterm{prop:g3:sorting-living-things:groups}{pez} de la comida que la charca produce''), la
        afirmación está mal.
\end{enumerate}
\end{method}

\begin{example}[La auditoría en marcha]\label{ex:g6:matter-from-food:audituse}
Audita el mirlo (\cref{ex:g2:what-animals-eat:winter}): entradas ---
lombrices, \omterm{prop:g3:sorting-living-things:groups}{insectos}, bayas; construido --- su propio crecimiento, las
\omterm{def:g1:animals-around-us:covering}{plumas} nuevas de cada año, la nidada de \omterm{def:g2:animals-grow-up:born}{huevos}; salidas --- excrementos
en el muro y la gran parte quemada que dio fuerza a un año de vuelos y
de cantos. El peso diario de comida de un pájaro cantor es enorme para
su tamaño justamente porque el vuelo quema muchísimo de la contabilidad.
\end{example}

\section{Ejercicios}

\begin{exercise}[$\star$]\label{exo:g6:matter-from-food:1}
Enuncia el oficio de los \omterm{def:g5:food-webs:roles}{consumidores} y contrástalo con el de las
\omterm{def:g5:food-webs:roles}{productoras}, en una frase cada uno.
\end{exercise}

\begin{exercise}[$\star$]\label{exo:g6:matter-from-food:2}
Enumera los pasos que convierten la hierba en \omterm{def:g3:bones-and-muscles:muscle}{músculo} de vaca,
nombrando el capítulo que trató cada uno.
\end{exercise}

\begin{exercise}[$\star$]\label{exo:g6:matter-from-food:3}
Da cuatro formas de \omterm{def:g6:matter-from-food:production}{producción de materia} de un \omterm{def:g1:animals-around-us:animal}{animal}, con un ejemplo
de cada una.
\end{exercise}

\begin{exercise}[$\star$]\label{exo:g6:matter-from-food:4}
Nombra los tres destinos del alimento de un \omterm{def:g5:food-webs:roles}{consumidor}, según la
\cref{prop:g6:matter-from-food:losses}.
\end{exercise}

\begin{exercise}[$\star$]\label{exo:g6:matter-from-food:5}
¿Por qué necesita una vaca lechera pienso todos los días para seguir
dando \omterm{def:g2:animals-grow-up:born}{leche}?
\end{exercise}

\begin{exercise}[$\star$]\label{exo:g6:matter-from-food:6}
Un \omterm{def:g2:animals-grow-up:born}{huevo} contiene \qty{60}{g} de \omterm{def:g6:plants-are-producers:matter}{materia orgánica}. ¿Dónde se fabricó
--- en la gallina o antes que ella? Desenrédalo con cuidado.
\end{exercise}

\begin{exercise}[$\star\star$]\label{exo:g6:matter-from-food:7}
Usa la \cref{prop:g6:matter-from-food:losses} para explicar por qué la
pirámide alimentaria se estrecha en cada piso.
\end{exercise}

\begin{exercise}[$\star\star$]\label{exo:g6:matter-from-food:8}
Pásale el \cref{met:g6:matter-from-food:audit} a un gato de casa:
entradas, construido y salidas --- con al menos dos apuntes por línea.
\end{exercise}

\begin{exercise}[$\star\star$]\label{exo:g6:matter-from-food:9}
``El zorro y la vaca hacen cosas distintas con la comida''. Corrige la
frase con la \cref{prop:g6:matter-from-food:menus} --- ¿qué es distinto
y qué es idéntico?
\end{exercise}

\begin{exercise}[$\star\star$]\label{exo:g6:matter-from-food:10}
Un niño crece \qty{3}{kg} en un año mientras se come varios cientos de
kilos de comida. ¿Adónde fue el resto? Responde con la contabilidad.
\end{exercise}

\begin{exercise}[$\star\star$]\label{exo:g6:matter-from-food:11}
La \omterm{def:g1:animals-around-us:covering}{concha} del caracol es sustancia mineral y, sin embargo, el
\cref{ex:g6:plants-are-producers:sorting} la contó como fabricación del
caracol. Reconcilia las dos cosas: ¿qué hizo el cuerpo del caracol, y
con qué suministros, para construirla?
\end{exercise}

\begin{exercise}[$\star\star\star$]\label{exo:g6:matter-from-food:12}
Una hectárea de terreno puede alimentar a más personas con patatas que
con el cerdo criado con esas patatas. Explícalo con la
\cref{prop:g6:matter-from-food:losses} --- ¿adónde va la diferencia
cuando las patatas dan el rodeo por el cerdo?
\end{exercise}

\section{Problema: La contabilidad de la granja}

\begin{problem}[{Problema de fin de semana --- una granja pequeña
auditada, del campo a la mesa}]\label{pb:g6:matter-from-food:1}
La pequeña finca de la abuela: un pasto con una vaca, un gallinero de
diez gallinas alimentadas con grano, una huerta y la cocina de la
familia. Este fin de semana llevas tú su contabilidad de materia.

\medskip\noindent\textbf{Parte I --- El piso de las \omterm{def:g5:food-webs:roles}{productoras}.}
\begin{enumerate}
  \item Enumera las \omterm{def:g5:food-webs:roles}{productoras} de la finca. ¿Cuáles son sus entradas
        minerales y cuál su fuente de \omterm{prop:g3:balanced-diet:jobs}{energía}?
  \item El pasto produce unos \qty{10000}{kg} de materia de hierba en
        una temporada. ¿En las \omterm{def:g1:plants-around-us:parts}{hojas} de quién se fabricó cada gramo?
  \item La abuela esparce el estiércol del gallinero por la huerta.
        ¿Qué regla de la \cref{prop:g5:people-and-nature:lasting} está
        cumpliendo y qué repone exactamente el estiércol?
  \item ¿Por qué no se le puede dar a la vaca directamente \omterm{prop:g4:what-plants-need:needs}{sales
        minerales} y agua, y ahorrarse el pasto?
\end{enumerate}

\medskip\noindent\textbf{Parte II --- El piso de los \omterm{def:g5:food-webs:roles}{consumidores}.}
\begin{enumerate}[resume]
  \item La vaca come alrededor de \qty{70}{kg} de hierba al día y da
        unos \qty{20}{kg} de \omterm{def:g2:animals-grow-up:born}{leche}. Nombra los tres destinos de los
        \qty{50}{kg} que faltan.
  \item Las gallinas convierten grano en \omterm{def:g2:animals-grow-up:born}{huevos}. Escribe la
        contabilidad de una gallina --- entradas, construido, salidas
        --- y di por qué una gallina en muda (que renueva todas sus
        \omterm{def:g1:animals-around-us:covering}{plumas}) pone menos \omterm{def:g2:animals-grow-up:born}{huevos} durante un tiempo.
  \item Un ratonero se lleva un pollito. Alarga hasta él la \omterm{def:g3:food-chains:chain}{cadena
        alimentaria} de la finca y di en qué eslabones se \emph{fabricó}
        materia y en cuáles solo se transformó.
  \item La familia come \omterm{def:g2:animals-grow-up:born}{huevos}, \omterm{def:g2:animals-grow-up:born}{leche} y verduras. Rastrea la \omterm{def:g6:plants-are-producers:matter}{materia
        orgánica} de un desayuno --- pan, \omterm{def:g2:animals-grow-up:born}{huevo} cocido, \omterm{def:g2:animals-grow-up:born}{leche} --- hasta
        la \omterm{def:g1:plants-around-us:parts}{hoja} donde se fabricó por primera vez.
\end{enumerate}

\medskip\noindent\textbf{Parte III --- La aritmética de los pisos.}
\begin{enumerate}[resume]
  \item La abuela dice: ``el pasto alimenta a una vaca; no alimentaría
        más que a una madriguera pequeña de zorros si los zorros
        comieran en él conejos criados con hierba''. Justifica su
        estimación con las pérdidas de cada piso.
  \item La huerta alimenta directamente a la familia. Con la lógica del
        \cref{exo:g6:matter-from-food:12}, explica por qué la huerta
        alimenta a más personas por metro cuadrado que la ruta
        pasto-vaca-leche.
  \item Con la sequía, la hierba deja de crecer. Predice, en orden, qué
        les pasa a la \omterm{def:g2:animals-grow-up:born}{leche} de la vaca, a su peso y a la línea de
        ``construido'' de la contabilidad --- y por qué no ayuda nada
        de sol de más mientras falle la lluvia (la clave está en la
        pregunta 7 del \cref{pb:g6:plants-are-producers:1}).
  \item Cierra la contabilidad con una frase por piso: \omterm{def:g5:food-webs:roles}{productoras} y
        \omterm{def:g5:food-webs:roles}{consumidores} --- y nombra a quien todavía falta en las cuentas
        (el tema del capítulo siguiente).
\end{enumerate}
\end{problem}
